\documentclass[a4paper]{article}
\usepackage{graphicx}
\usepackage{amsmath,amssymb}
\usepackage{hyperref}
\usepackage{minted}
\usepackage{multirow}

\title{awarness_vs_conciousness.tex}
\date{}

\begin{document}
    \maketitle
    
    
    \section{Terminology and Scope}
        
        In everyday English, \emph{awareness} and \emph{consciousness} are often used interchangeably, but in philosophy of mind, cognitive science, and neuroscience they typically refer to different levels of description. This section fixes working definitions and clarifies how the terms relate.
    
    
    \section{Awareness}
        
        \emph{Awareness} refers to being aware \emph{of something}. It is usually object-directed: one is aware of a sound, a pain, a visual object, a thought, or a change in the environment. In many scientific contexts, awareness is treated as a capacity for detection, monitoring, or access to information, and it can be discussed without committing to any specific theory about subjective experience.
    
    
    \section{Consciousness}
        
        \emph{Consciousness} is the broader notion. At a minimum, it can mean being in a conscious state (as opposed to being under deep anesthesia or in a coma). In many theoretical discussions, consciousness also includes the presence of \emph{subjective experience} (often described as there being something it is like to be in that state). Some usages further distinguish forms such as phenomenal consciousness (experience), access consciousness (availability of information for report and control), and self-consciousness (awareness of oneself as oneself).
    
    
    \section{Relationship Between the Terms}
        
        A common working relationship is:
        \begin{itemize}
            \item Awareness is typically treated as a component or manifestation within a conscious state.
            \item Consciousness is not reducible to awareness alone, because it can include the global condition of being conscious and, in many accounts, the qualitative character of experience.
        \end{itemize}
        
        Illustrative cases:
        \begin{itemize}
            \item Deep anesthesia: neither consciousness nor awareness is present.
            \item Dreaming sleep: consciousness (as experience) may be present, while awareness of the external environment is reduced or absent.
            \item Normal wakefulness: both consciousness and multiple forms of awareness (of environment, body, thoughts) can be present.
        \end{itemize}
    
    
    \section{A Compact Distinction}
        
        A concise way to separate the ideas is:
        \begin{itemize}
            \item Awareness: access or sensitivity to particular contents (being aware \emph{of} X).
            \item Consciousness: the overall condition in which experience occurs (and, depending on the theory, the subjective character of that experience).
        \end{itemize}
    % \bibliographystyle{plain}
    % \bibliography{/home/donkarlo/Dropbox/repo/shared_research/refs/refs}
\end{document}